\documentclass[11pt,a4paper]{article}
\usepackage[utf8]{inputenc}
\usepackage[T1]{fontenc}
\usepackage[french]{babel}
\usepackage[top=2cm,bottom=2cm,left=2cm,right=2cm]{geometry}
\usepackage{xcolor}
\usepackage{tcolorbox}
\tcbuselibrary{skins,breakable}
\usepackage{fancyhdr}
\usepackage{enumitem}
\usepackage{lmodern}

% --- Couleur discipline : ÉCONOMIE ---
\definecolor{mainColor}{RGB}{0,80,160}
\definecolor{bgLight}{RGB}{245,248,255}
\definecolor{darkGray}{RGB}{50,50,50}

% --- Boîtes ---
\tcbset{boxrule=0.5pt, arc=3pt, left=5pt, right=5pt, top=3pt, bottom=3pt, breakable}

\newtcolorbox{defbox}[1]{
  colback=mainColor!8, colframe=mainColor, coltitle=white,
  fonttitle=\bfseries\footnotesize, title=DÉFINITION \textemdash\ #1, left=5pt}

\newtcolorbox{keybox}[1][]{
  colback=darkGray!7, colframe=darkGray!50, coltitle=white,
  fonttitle=\bfseries\footnotesize, title=POINT CLÉ, left=5pt}

\newtcolorbox{exbox}{
  colback=mainColor!5, colframe=mainColor!40,
  fonttitle=\bfseries\footnotesize, title=EXEMPLE, left=5pt}

% --- En-tête ---
\pagestyle{fancy}\fancyhf{}
\fancyhead[L]{\small\textcolor{mainColor}{\textbf{CEJM — BTS SIO}}}
\fancyhead[C]{\small\textbf{Chapitre 1 — Les agents économiques}}
\fancyhead[R]{\small\textcolor{mainColor}{\textbf{ÉCONOMIE}}}
\fancyfoot[C]{\small\thepage}
\renewcommand{\headrulewidth}{0.5pt}
\setlength{\headheight}{15pt}
\setlist{itemsep=2pt, topsep=3pt, parsep=0pt}

\begin{document}

\begin{center}
  {\Large\bfseries\textcolor{mainColor}{Chapitre 1}}\\[2pt]
  {\large Les agents économiques en relation avec l'entreprise}\\[4pt]
  \textit{\textcolor{darkGray}{Comment l'entreprise interagit-elle avec les autres agents économiques ?}}
\end{center}
\vspace{4pt}\hrule\vspace{8pt}

% =========================================================
\section*{Définitions essentielles}

\begin{defbox}{Agent économique}
Individu ou groupe d'individus constituant un \textbf{centre de décision économique indépendant}, regroupé selon sa fonction principale dans l'économie.
\end{defbox}

\vspace{4pt}
\begin{defbox}{Production marchande}
Biens et services vendus à un \textbf{prix de marché} (> 50\,\% du coût de production). Assurée par les \textbf{entreprises}.
\end{defbox}

\vspace{4pt}
\begin{defbox}{Production non marchande}
Biens et services fournis \textbf{gratuitement ou quasi-gratuitement} (< 50\,\% du coût). Assurée par les \textbf{administrations publiques}.
\end{defbox}

\vspace{4pt}
\begin{defbox}{Circuit économique}
Représentation schématique des \textbf{flux} (réels et monétaires) échangés entre les agents économiques.
\end{defbox}

% =========================================================
\section*{Les 4 agents économiques principaux}

\subsection*{1. Les ménages}
\begin{keybox}
Fonction principale : \textbf{CONSOMMATION} de biens et services pour satisfaire leurs besoins.
\end{keybox}
\begin{itemize}
  \item Offrent leur \textbf{travail} sur le marché du travail (→ salaires)
  \item Consomment et \textbf{épargnent} une partie de leur revenu
  \item Paient des \textbf{impôts et cotisations} à l'État
\end{itemize}

\subsection*{2. Les entreprises}
\begin{keybox}
Fonction principale : \textbf{PRODUCTION} de biens et services marchands pour réaliser un profit.
\end{keybox}
\begin{itemize}
  \item Demandeurs de \textbf{travail} (ménages) et de \textbf{capitaux} (banques)
  \item Relations B2B (fournisseurs/clients entreprises)
  \item Versent impôts et cotisations sociales
\end{itemize}

\subsection*{3. Les banques (institutions financières)}
\begin{keybox}
Fonction principale : \textbf{FINANCEMENT} de l'économie par collecte de l'épargne et octroi de crédits.
\end{keybox}
\begin{itemize}
  \item Collectent l'épargne des ménages
  \item Accordent des crédits aux entreprises et ménages
  \item Le prix du capital = \textbf{taux d'intérêt}
\end{itemize}

\subsection*{4. Les administrations publiques (APU)}
\begin{keybox}
Fonction principale : Produire des \textbf{services non marchands} et répondre aux besoins d'intérêt général.
\end{keybox}
\begin{itemize}
  \item État, collectivités locales, Sécurité sociale
  \item Financées par \textbf{impôts, taxes, cotisations sociales}
  \item Versent des \textbf{prestations sociales} aux ménages
\end{itemize}

% =========================================================
\section*{Les marchés qui relient les agents}

\begin{center}
\begin{tabular}{|l|l|l|}
\hline
\textbf{Marché} & \textbf{Flux réel} & \textbf{Flux monétaire} \\
\hline
Biens et services & Produits/services & Paiements (CA) \\
Travail & Force de travail & Salaires \\
Capitaux & Épargne/crédits & Intérêts \\
Financier & Titres (actions, obligations) & Dividendes \\
\hline
\end{tabular}
\end{center}

\vspace{4pt}
\begin{exbox}
\textbf{Bpifrance} : banque publique qui finance les projets d'entreprises (innovation, export, fonds propres) — lien entre APU et entreprises.
\end{exbox}

% =========================================================
\section*{Points clés à retenir}

\begin{keybox}
L'entreprise est au centre d'un réseau de relations avec \textbf{tous} les agents économiques.
\end{keybox}
\vspace{3pt}
\begin{keybox}
Les échanges génèrent 2 types de flux : \textbf{réels} (biens, services, travail) et \textbf{monétaires} (prix, salaires, intérêts).
\end{keybox}
\vspace{3pt}
\begin{keybox}
La \textbf{conjoncture économique} (croissance, chômage, inflation) influence les comportements de tous les agents.
\end{keybox}

% =========================================================
\section*{Mots-clés}
\textcolor{mainColor}{\textbf{Agent économique}} $\cdot$
\textcolor{mainColor}{\textbf{Production marchande / non marchande}} $\cdot$
\textcolor{mainColor}{\textbf{Circuit économique}} $\cdot$
\textcolor{mainColor}{\textbf{Flux réels / monétaires}} $\cdot$
\textcolor{mainColor}{\textbf{Marché du travail}} $\cdot$
\textcolor{mainColor}{\textbf{Marché financier}} $\cdot$
\textcolor{mainColor}{\textbf{Taux d'intérêt}} $\cdot$
\textcolor{mainColor}{\textbf{Conjoncture économique}} $\cdot$
\textcolor{mainColor}{\textbf{Financement direct / indirect}}

\end{document}
